Primary screening must expose enough GPU work without losing a candidate
that should reach residual refinement. A query-centric scan launches one
block per query; a small batch therefore offers few blocks even when each
query visits many vectors. GPU selection is a substantial design problem
in its own right \cite{johnson_gpu}. The grouped-scan component, introduced
in v62, partitions the visited candidates across blocks and combines their
survivors. Its contract is exact selection under the computed primary
scores. IVF routing and the quantised representation remain approximate.

\subsection{Candidate partition and deterministic ordering}

For a fixed query, concatenate the selected inverted lists in probe order
to obtain $n_c$ candidate occurrences. Assume that each visited database
row occurs once and that the index stays fixed during search. For
$G=\gscan\in\{1,2,4,8\}$, block $(q,g)$ owns the interval
\begin{equation}
 U_g=\left[\left\lfloor\frac{g n_c}{G}\right\rfloor,
              \left\lfloor\frac{(g+1)n_c}{G}\right\rfloor\right),
 \qquad g=0,\ldots,G-1.
 \label{eq:5-partition}
\end{equation}
Here an interval denotes the corresponding candidates, rather than their
database IDs. The device computes its endpoints from the routed candidate
count. Groups differ in size by at most one candidate; a long list can
cross a group boundary. This balances candidate counts, whereas assigning
equal numbers of lists need not do so. Each thread finds its first list
through the probe-prefix offsets and advances its list cursor while
scanning its interval. Residual codes are not read at this stage.

The selection key is
\begin{equation}
 K_q(x)=\bigl(\primarydist(q,x),p(x)\bigr),
 \qquad
 (a,p)\prec(b,r)\iff a<b\ \text{or}\ (a=b\ \text{and}\ p<r),
 \label{eq:5-primary-key}
\end{equation}
where $p(x)$ is the physical row position in the fixed inverted-list
storage. It is unique among the query's candidates and independent of
the order of atomic appends. It is not the database ID used to break
final-result ties. With non-NaN scores, this defines a total order on
distinct candidates. All local comparisons and merge decisions use that
same order. Partition-independent selection additionally requires that
each candidate receive the same computed score in every compared path;
the ordering rule alone cannot remove floating-point scoring differences.

\subsection{Exact local retention and global coverage}

Write $\operatorname{Top}_c(V)$ for the first $\min(c,|V|)$ candidates of
$V$ sorted by $K_q$, with $c=\max\{k,\lceil\alpha k\rceil\}$. Each
independent group retains its full local top-$c$, not a share $c/G$.
This distinction is necessary: with $c=4$ and two groups, all four
global winners may belong to one group. Retaining only two from each
group loses winners even though the combined capacity is four.

\begin{proposition}[Local coverage and truncated merging]
\label{prop:5-local-coverage}
For disjoint groups $U_0,\ldots,U_{G-1}$, fixed candidate scores and a
common total order,
\begin{equation}
 \operatorname{Top}_c\!\left(\bigcup_g U_g\right)
 =\operatorname{Top}_c\!\left(\bigcup_g
                    \operatorname{Top}_c(U_g)\right).
 \label{eq:5-topc-union}
\end{equation}
Pairwise sorted merging with truncation to $c$ after every merge preserves
this result, including when fewer than $c$ candidates exist.
\end{proposition}
\begin{proof}
Let $x$ belong to the global top-$c$ and to group $U_g$. Fewer than $c$
candidates precede $x$ globally, so fewer than $c$ precede it within $U_g$.
Consequently $x\in\operatorname{Top}_c(U_g)$, and the local union contains
every global winner. Selecting its first $c$ therefore recovers precisely
the global winners in the same order. If the whole universe has fewer
than $c$ members, every member is retained.

For merging, the first $c$ of two sorted lists are their union's
top-$c$. A discarded candidate already has at least $c$ predecessors
within that union and cannot enter the top-$c$ of any larger union.
Induction over the merge tree establishes the claim.
\end{proof}

\begin{figure}[tbp]
\centering
% Schematic only: illustrative keys, not experimental observations.
\begin{tikzpicture}[
  font=\small\linespread{1}\selectfont,
  box/.style={draw=ThesisNavy!65,rounded corners=2pt,align=center,
    inner sep=5pt,text width=2.65cm,minimum height=2.05cm},
  merged/.style={box,text width=5.55cm,minimum height=1.05cm,
    fill=ThesisNavy!5},
  arr/.style={-{Latex[length=2mm]},draw=ThesisNavy!80,thick}]
\node[box] (g0) at (0,0) {$\operatorname{Top}_3(U_0)$\\[2pt]
  $\mathbf{(1,8)}$\\$\mathbf{(2,6)}$\\$(2,9)$};
\node[box] (g1) at (3.5,0) {$\operatorname{Top}_3(U_1)$\\[2pt]
  $(4,1)$\\$(5,4)$\\\textit{length 2}};
\node[box] (g2) at (7,0) {$\operatorname{Top}_3(U_2)$\\[2pt]
  $\varnothing$\\\textit{length 0}};
\node[box] (g3) at (10.5,0) {$\operatorname{Top}_3(U_3)$\\[2pt]
  $\mathbf{(2,3)}$\\$(7,2)$\\$(8,5)$};
\node[merged] (m0) at (1.75,-2.5) {Merge and retain first 3\\[2pt]
  $\mathbf{(1,8)},\ \mathbf{(2,6)},\ (2,9)$};
\node[merged] (m1) at (8.75,-2.5) {Merge and retain first 3\\[2pt]
  $\mathbf{(2,3)},\ (7,2),\ (8,5)$};
\node[merged,text width=6.8cm,fill=ThesisNavy!12] (out) at (5.25,-4.3)
  {Global primary top-3\\[2pt]
   $\mathbf{(1,8)},\ \mathbf{(2,3)},\ \mathbf{(2,6)}$};
\draw[arr] (g0.south) -- (m0.north west);
\draw[arr] (g1.south) -- (m0.north east);
\draw[arr] (g2.south) -- (m1.north west);
\draw[arr] (g3.south) -- (m1.north east);
\draw[arr] (m0.south) -- (out.north west);
\draw[arr] (m1.south) -- (out.north east);
\end{tikzpicture}

\caption{Local coverage and truncated pairwise merging, illustrated with
$c=3$ and $G=4$. Entries are $(\primarydist,p)$ keys, not measured data.
Bold entries are the final global winners. A short group contributes its
valid prefix; an empty group contributes no entries. Every merge uses the
same key order and retains at most $c$ entries.}
\label{fig:5-global-selection}
\end{figure}

The CUDA implementation maintains a shared-memory buffer of $C$ key
pairs, where $C$ is a power of two and $C\geq c+T$ for a scan block of
$T$ threads. A compaction sorts this buffer with a bitonic network,
retains at most $c$ valid entries and records the last retained key as
the local threshold when $c$ entries exist. Before that point the
threshold is a sentinel above every valid key. Between compactions,
only candidates preceding the threshold are appended. Its $c$ witness
candidates make the threshold a safe upper bound for the best $c$ of
the scanned prefix, even as further candidates arrive.

Each chunk contributes at most $T$ appends. After a chunk, the block
compacts whenever the next chunk would exceed capacity. Following
compaction at most $c$ entries remain, so $C\geq c+T$ provides space
for the next chunk; otherwise the explicit capacity check provides it.
This invariant prevents overflow without restricting how many competitive
candidates one thread may encounter. A per-thread fixed-size retention
buffer does not provide the same coverage guarantee.

An alternative is to materialise all $n_c$ score/position pairs and then
invoke a separate selector. That preserves the target but adds an
$O(n_c)$ intermediate array and its writes. Here packed-code scoring
feeds a bounded local buffer directly, exchanging repeated local sorting
for $O(Gc)$ intermediate output. Register-based selection is another
established GPU design \cite{johnson_gpu}; this campaign evaluates the
chosen scan variants rather than demonstrating superiority over every
library selector.

\subsection{Merge implementation, storage and boundary cases}

For $G>1$, one 256-thread block per query performs $\log_2 G$ rounds of
pairwise merging. A pair with valid lengths $a,b$ emits
$\min(c,a+b)$ keys. Merge Path partitions this output into contiguous
thread slices: a binary search finds the input offsets for each slice,
then the thread merges its slice sequentially \cite{gpu_merge_path}.
All comparisons use \cref{eq:5-primary-key}; block barriers separate
rounds. \Cref{fig:5-global-selection} illustrates the truncation invariant.
The tree emits at most $c(G-1)$ key pairs across its internal nodes,
in addition to the partition searches and synchronisation work.

The two global score/position workspaces have capacities
$[b_{\rm cap},G,c]$ and $[b_{\rm cap},G/2,c]$ and alternate between
rounds. With four-byte scores and positions, their combined payload is
\begin{equation}
 W_{\rm pairs}=8b_{\rm cap}c\left(G+\frac{G}{2}\right)
 \quad\text{bytes},\qquad G>1.
 \label{eq:5-group-workspace}
\end{equation}
Group counts require another $4b_{\rm cap}G$ bytes; shared bounds add
$8b_{\rm cap}$ when provisioned. These are additional to final survivor
arrays and other search workspace. Allocation grows before graph capture
and is retained for reuse, so choosing a smaller $G$ later does not
necessarily release the reserved memory. The merge keeps only list
lengths in shared memory rather than all $Gc$ keys.

Counts, rather than sentinel scores alone, delimit local valid prefixes.
Compaction pads its sorting buffer with $(+\infty,\mathrm{INT\_MAX})$;
valid positions must be below that reserved position. Empty groups have
count zero. The final survivor array pads unused slots with
$(+\infty,-1)$, and refinement skips negative positions. These two
sentinels serve different purposes: the negative padding position must
not participate as a valid key in a merge. $G=1$ writes directly to the
survivor arrays and omits the merge.

The older query-centric path is selected by the implementation label
$G=0$, meaning one block per query, not zero blocks. Its default primary
selection is exact by score, but its distance-only cutoff can retain a
different member of a tie from the complete-key path. The lossy
\texttt{K\_LOCAL=4} alternative requires
\texttt{JHQ\_EXACT\_TOPCK=0} in a build that permits it; paper-faithful
builds force exact selection. It is not the measured default.
\Cref{sec:7-4} reports 64 grouped-scan comparisons with no non-tie primary
differences and identical final outputs on the used queries. Those checks
support the implementation within their coverage. The proposition
establishes a conditional selection property, not raw-space exact kNN or
unconditional equivalence to every legacy tie policy.
